Die Spannenergie der Feder ist mit $k = F/s$
\begin{align}
 E_\mathrm{el} &= \frac{1}{2}k\,s^2 = \EsfF = \frac{1}{2}\cdot\Fsf\cdot\ssf = \Esf
\end{align}
Wird diese vollständig in potentielle Energie umgewandelt, $E_\mathrm{el} = m\,g\,h$, so steigt die Feder um
\begin{align}
 h &= \hsfF \\
   &= \frac{\Fsf\cdot\ssf}{2\cdot\msf\cdot\gE} \\
   &= \hsf \approx \result{\hsfP{-2}{2}}
\end{align}
Gemeint ist die Höhe, um die der Schwerpunkt der Feder steigt. Reibung, Luftwiderstand und die Schwingungsenergie der Feder nach dem Absprung sind vernachlässigt.
